\clearpage\thispagestyle{empty}
\pdfbookmark[0]{Quellen zum Pallas-Bericht vom Juni 1802}{pallas180206-sources}
\begin{center}\textbf{Quellen zum Pallas-Bericht vom Juni 1802}\end{center}
\noindent Diese KI-Übertragung in heutiges Deutsch folgt dem deutschen Auszug in Carl Friedrich Gauß’ \href{https://gdz.sub.uni-goettingen.de/id/PPN236007467}{\textit{Werken}, Band~VI} (Göttingen, 1874), S.~212--216. Franz Xaver von Zach berichtet über die Entdeckung der Pallas durch Heinrich Wilhelm Matthias Olbers. Der Auszug enthält Rechnungen von Carl Friedrich Gauß und zitiert aus einem seiner Briefe.
\par\smallskip\noindent Von Zachs \href{https://api.digitale-sammlungen.de/iiif/image/v2/bsb10538597_00607/full/1450,/0/default.jpg}{ursprünglicher Aufsatz}, Nr.~LX, erschien im Juniheft 1802 der Zeitschrift \textit{Monatliche Correspondenz zur Beförderung der Erd- und Himmels-Kunde} (Gotha, Becker), Band~V, S.~591--606. Der Auszug folgt S.~591--600 und setzt mit der \href{https://api.digitale-sammlungen.de/iiif/image/v2/bsb10538597_00624/full/1450,/0/default.jpg}{Ephemeride auf S.~606} wieder ein. Der Verweis auf den erfolglosen Ansatz einer Kreisbahn führt ins \href{https://api.digitale-sammlungen.de/iiif/image/v2/bsb10538597_00498/full/1450,/0/default.jpg}{vorige Heft, S.~486}. Die Vergleiche der ersten und zweiten Bahnelemente mit Beobachtungen stehen auf \href{https://api.digitale-sammlungen.de/iiif/image/v2/bsb10538597_00610/full/1450,/0/default.jpg}{S.~594} und \href{https://api.digitale-sammlungen.de/iiif/image/v2/bsb10538597_00614/full/1450,/0/default.jpg}{S.~596}.
\par\smallskip\noindent Beim ersten Vergleich steht für Olbers’ Beobachtung vom 1. Mai in den \textit{Werken} der Unterschied der Rektaszension \(+20.0\). Im Zeitschriftendruck steht \(+5.0\). Der später angegebene Bahnbogen beträgt in den \textit{Werken} \(7^\circ 26'\), im \href{https://api.digitale-sammlungen.de/iiif/image/v2/bsb10538597_00611/full/1450,/0/default.jpg}{Zeitschriftendruck} \(7^\circ 23'\). Die Übertragung folgt den Zahlen der \textit{Werke} und bewahrt auf S.~214 deren Wortlaut \textit{einfachen} in der Aussage zur Parallaxenrechnung. Die Auslassung vor der Ephemeride folgt dem Auszug nach dem Absatz auf \href{https://api.digitale-sammlungen.de/iiif/image/v2/bsb10538597_00618/full/1450,/0/default.jpg}{S.~600}.
\par\smallskip\noindent Übertragung und Satz wurden mit KI erstellt. Die Digitalisate der Zeitschrift stammen von der Bayerischen Staatsbibliothek.
\clearpage
